% This is samplepaper.tex, a sample chapter demonstrating the
% LLNCS macro package for Springer Computer Science proceedings;
% Version 2.20 of 2017/10/04
%
\documentclass[runningheads]{llncs}
%
\usepackage{graphicx}
\usepackage{amsmath}
\usepackage{amssymb}
\usepackage{booktabs}
\usepackage{multirow}

\newcommand\red[1]{\color{red}#1}
% Used for displaying a sample figure. If possible, figure files should
% be included in EPS format.
%
% If you use the hyperref package, please uncomment the following line
% to display URLs in blue roman font according to Springer's eBook style:
% \renewcommand\UrlFont{\color{blue}\rmfamily}

\begin{document}
%
% \title{AME-CAM: Attentive Multiple-Exit CAM for Weakly Supervised Segmentation on MRI Brain Tumor}
%
%\titlerunning{Abbreviated paper title}
% If the paper title is too long for the running head, you can set
% an abbreviated paper title here
%
% \author{First Author\inst{1}\orcidID{0000-1111-2222-3333} \and
% Second Author\inst{2,3}\orcidID{1111-2222-3333-4444} \and
% Third Author\inst{3}\orcidID{2222--3333-4444-5555}}
% %
% \authorrunning{F. Author et al.}
% % First names are abbreviated in the running head.
% % If there are more than two authors, 'et al.' is used.
% %
% \institute{Princeton University, Princeton NJ 08544, USA \and
% Springer Heidelberg, Tiergartenstr. 17, 69121 Heidelberg, Germany
% \email{lncs@springer.com}\\
% \url{http://www.springer.com/gp/computer-science/lncs} \and
% ABC Institute, Rupert-Karls-University Heidelberg, Heidelberg, Germany\\
% \email{\{abc,lncs\}@uni-heidelberg.de}}
%
% \maketitle              % typeset the header of the contribution

% \section{Comparison on the TCGA-LGG dataset}
\begin{table}
\centering
\caption{Comparison with weakly supervised methods (WSSS), unsupervised method (UL), and fully supervised methods (FSL) on The Cancer Genome Atlas Lower-Grade Glioma collection (TCGA-LGG) dataset. Results are reported in the form of mean$\pm$std. We mark the highest score among WSSS methods with bold text.}
\label{Table:Supple_TCGA-LGG}
\begin{tabular}{c|c|c|c|l}
\hline
Type                  & Method            & Dice $\uparrow$       & IoU $\uparrow$        & \multicolumn{1}{c}{HD95 $\downarrow$} \\ \hline
\multirow{6}{*}{WSSS} & Grad-CAM (2016)   & 0.407$\pm$0.229          & 0.283$\pm$0.190          & 81.278$\pm$38.648                        \\
                      & ScoreCAM (2020)   & 0.490$\pm$0.222          & 0.351$\pm$0.184          & 57.821$\pm$30.275                        \\
                      & LFI-CAM (2021)    & 0.451$\pm$0.242          & 0.323$\pm$0.206          & 62.077$\pm$36.067                        \\
                      & LayerCAM (2021)   & 0.378$\pm$0.211          & 0.255$\pm$0.171          & 69.696$\pm$42.493                        \\
                      & Swin-MIL (2022)   & 0.271$\pm$0.144          & 0.165$\pm$0.101          & 111.634$\pm$17.670                       \\
                      & AME-CAM (ours)    & \textbf{0.517$\pm$0.244} & \textbf{0.385$\pm$0.227} & \textbf{64.766$\pm$35.706}               \\ \hline
UL                    & C\&F (2020)       & 0.300$\pm$0.162          & 0.188$\pm$0.119          & 101.980$\pm$18.358                       \\ \hline
\multirow{2}{*}{FSL}  & C\&F (2020)       & 0.637$\pm$0.216          & 0.502$\pm$0.224          & 46.925$\pm$17.889                        \\
                      & Opt. U-Net (2021) & 0.668$\pm$0.196          & 0.532$\pm$0.209          & 47.035$\pm$19.529                        \\ \hline 
\end{tabular}
\end{table}

% The Cancer Genome Atlas Lower-Grade Glioma collection (TCGA-LGG) dataset \cite{pedano2016radiology} consists of 110 patients with 3,600 2D images collected from the TCGA lower-grade glioma collection who had preoperative imaging data available, containing at least a fluid-attenuated inversion recovery (FLAIR) sequence. We follow the preprocess step and the split rate in the BraTS dataset to separate the dataset into training, validation, and testing (180 positive and 351 negative images) for experiments.


% \section{Ablation study of Using Different Loss Function in Activation Aggregation}
\begin{table}[h]
\centering
% \setlength{\tabcolsep}{3pt}
\caption{Ablation study for using different loss function in the aggregation phase, conventional cross-entropy loss v.s. the proposed C$^2$AM loss. The experiments are done with the BraTS dataset. Results are reported in the form of mean$\pm$std. We mark the highest score among WSSS methods with bold text.}
\label{Table:Supple_agg_phase_loss}
\begin{tabular}{c|c|c|c|l}
\hline
Dataset                                       & Method         & Dice $\uparrow$       & IoU $\uparrow$        & \multicolumn{1}{c}{HD95 $\downarrow$} \\ \hline
\multirow{2}{*}{\shortstack{BraTS\\T1}}       & Cross-Entropy  & 0.494$\pm$0.174          & 0.345$\pm$0.154          & 43.903$\pm$34.428                        \\
                                              & C$^2$AM (ours) & \textbf{0.631$\pm$0.119} & \textbf{0.471$\pm$0.119} & \textbf{21.813$\pm$18.219}               \\ \hline
\multirow{2}{*}{\shortstack{BraTS\\T1-CE}}    & Cross-Entropy  & 0.644$\pm$0.142          & 0.489$\pm$0.145          & 20.255$\pm$12.525                        \\
                                              & C$^2$AM (ours) & \textbf{0.695$\pm$0.095} & \textbf{0.540$\pm$0.108} & \textbf{18.129$\pm$12.335}               \\ \hline
\multirow{2}{*}{\shortstack{BraTS\\T2}}       & Cross-Entropy  & 0.684$\pm$0.155          & 0.539$\pm$0.167          & 27.966$\pm$32.604                        \\
                                              & C$^2$AM (ours) & \textbf{0.721$\pm$0.086} & \textbf{0.571$\pm$0.101} & \textbf{14.940$\pm$8.736}                \\ \hline
\multirow{2}{*}{\shortstack{BraTS\\T2-FLAIR}} & Cross-Entropy  & 0.819$\pm$0.105          & 0.705$\pm$0.140          & 12.039$\pm$9.160                         \\
                                              & C$^2$AM (ours) & \textbf{0.862$\pm$0.088} & \textbf{0.767$\pm$0.122} & \textbf{8.664$\pm$6.440}                 \\ \hline
\end{tabular}
\end{table}
% In this section, we explore the ablation study of using different loss functions in the activation aggregation stage. The dice comparison of using conventional cross-entropy loss and the proposed C$^2$AM on the BraTS dataset is shown in Table \ref{Table:Supple_agg_phase_loss}. From the table, we can observe that the unsupervised contrastive loss C$^2$AM outperforms the cross-entropy loss in all metrics of all image modalities. The main difference is that C$^2$AM is a pixel-wise contrastive loss, while the cross-entropy loss requires a global average pooling operation, which somehow causes the information lost. Thus, we still recommend using C$^2$AM for training in the activation aggregation stage.

% \section{Full Ablation study of Using Single-Exit and Multiple-Exit}
% \input{Table/Supple_single_exit}
% To show that the use of multiple-exit in the activation extraction stage is reasonable, we compare the proposed AME-CAM with using four single-exit from the results of four internal classifiers $M_1$ to $M_4$ (referred to Fig. 1 in the main paper) in Table \ref{Table:ablation_single_exit}. The comparisons show that the activation map obtained from the shallow layer $M_1$ and the deepest layer $M_4$ result in low dice scores, which are around 0.2 in positive patients. This is because, in the shallow layer, the network is not deep enough to learn the tumor region yet. In addition, the resolution of the activation map obtained from the deepest layer is too low to contain sufficient information to make a clear boundary for the tumor.
% Results of the internal classifier from the middle of the network ($M_2$ and $M_3$) achieve the highest dice score and IoU, which are about 0.5 and 0.4 respectively in positive patients, and the dice score of over 0.7 in the overall evaluation.

% To evaluate whether using results from all internal classifiers will lead to the highest performance, we further apply the proposed method to the two internal classifiers of the highest dice, i.e., $M_2$ and $M_3$, called $M_2+M_3$. Compared with using all internal classifiers ($M_1$ to $M_4$), $M_2+M_3$ results in 16.9\% and 19.2\% lower dice and IoU of positive patients. However, AME-CAM also allows to aggregate activation from $M_1$ and $M_4$, which increase the possibility of generating false activation region. Compared with only using $M_2$ and $M_3$, it leads to 1.9\% lower dice score of negative patients. Considering the overall evaluation, though with a slightly lower dice in negative patients, the AME-CAM still has the optimal performance among all the experiments of single-exit and multiple-exit.

% \section{Comparison with WSSS Methods using ResNet-50 and the BraTS dataset}
\begin{table}
\centering
\caption{Comparison with WSSS Methods using ResNet-50 on BraTS dataset with T1, T1-CE, T2, and T2-FLAIR MRI images. Results are reported in the form of mean$\pm$std. We mark the highest score among WSSS methods with bold text.}
\label{Table:Supple_Res50}
% \setlength{\tabcolsep}{4pt}
\begin{tabular}{c|c|c|l}
\hline
\multicolumn{4}{c}{BraTS T1} \\ \hline
Method          & Dice $\uparrow$ & IoU $\uparrow$ & \multicolumn{1}{c}{HD95 $\downarrow$} \\ \hline
Grad-CAM (2016) & 0.046$\pm$0.086    & 0.026$\pm$0.053   & 132.469$\pm$33.056                       \\
ScoreCAM (2020) & 0.233$\pm$0.089    & 0.134$\pm$0.058   & 75.245$\pm$15.082                        \\
LFI-CAM (2021)  & 0.166$\pm$0.159    & 0.100$\pm$0.107   & 132.215$\pm$48.849                       \\
LayerCAM (2021) & 0.453$\pm$0.195    & 0.313$\pm$0.164   & 40.190$\pm$55.079                        \\
AME-CAM  (ours) & \textbf{0.587$\pm$0.149}    & \textbf{0.430$\pm$0.141}   & \textbf{20.903$\pm$19.565 }                       \\ \hline
\multicolumn{4}{c}{BraTS T1-CE} \\ \hline
Method          & Dice $\uparrow$ & IoU $\uparrow$ & \multicolumn{1}{c}{HD95 $\downarrow$} \\ \hline
Grad-CAM (2016) & 0.341$\pm$0.132    & 0.213$\pm$0.096   & 48.055$\pm$17.094                        \\
ScoreCAM (2020) & 0.343$\pm$0.235    & 0.234$\pm$0.188   & 50.868$\pm$32.598                        \\
LFI-CAM (2021)  & 0.639$\pm$0.136    & 0.484$\pm$0.143   & \textbf{18.686$\pm$8.465}                         \\
LayerCAM (2021) & 0.480$\pm$0.213    & 0.341$\pm$0.176   & 40.003$\pm$72.395                        \\
AME-CAM  (ours) & \textbf{0.652$\pm$0.099}    & \textbf{0.491$\pm$0.107}   & 19.327$\pm$10.951                        \\ \hline
\multicolumn{4}{c}{BraTS T2} \\ \hline
Method          & Dice $\uparrow$ & IoU $\uparrow$ & \multicolumn{1}{c}{HD95 $\downarrow$} \\ \hline
Grad-CAM (2016) & 0.075$\pm$0.060    & 0.040$\pm$0.035   & 147.272$\pm$26.227                       \\
ScoreCAM (2020) & 0.479$\pm$0.219    & 0.342$\pm$0.193   & 38.509$\pm$27.572                        \\
LFI-CAM (2021)  & 0.523$\pm$0.209    & 0.383$\pm$0.205   & 65.917$\pm$56.581                        \\
LayerCAM (2021) & 0.638$\pm$0.184    & 0.493$\pm$0.190   & 26.907$\pm$33.535                        \\
AME-CAM  (ours) & \textbf{0.682$\pm$0.168}    & \textbf{0.540$\pm$0.175}   & \textbf{19.078$\pm$24.685}                        \\ \hline
\multicolumn{4}{c}{BraTS T2-FLAIR} \\ \hline
Method          & Dice $\uparrow$ & IoU $\uparrow$ & \multicolumn{1}{c}{HD95 $\downarrow$} \\ \hline
Grad-CAM (2016) & 0.221$\pm$0.084    & 0.127$\pm$0.056   & 68.666$\pm$22.125                        \\
ScoreCAM (2020) & 0.588$\pm$0.252    & 0.458$\pm$0.239   & 26.126$\pm$23.682                        \\
LFI-CAM (2021)  & 0.739$\pm$0.164    & 0.611$\pm$0.197   & 14.429$\pm$6.620                         \\
LayerCAM (2021) & 0.638$\pm$0.196    & 0.497$\pm$0.200   & 23.290$\pm$35.905                        \\
AME-CAM  (ours) & \textbf{0.759$\pm$0.132}    & \textbf{0.629$\pm$0.163}   & \textbf{12.696$\pm$8.610}                         \\ \hline
\end{tabular}
\end{table}



% In this section, we show the quantitative comparison among the CAM-based WSSS methods using ResNet-50 as the classifier backbone in Table \ref{Table:Compare_CAM_Res50}. 
% In all modalities of positive patients, the dice scores of LayerCAM[9] are extremely low, which are lower than 0.3.
% As for the ScoreCAM[23] and the LFI-CAM[13], they achieve a significant improvement from LayerCAM with a higher dice score, except for the T1 modality of LFI-CAM. The proposed AME-CAM achieves the optimal dice score in the evaluation of positive patients.
% When it comes to the negative cases, all the previous works still suffer from low dice scores in negative patients, which is due to the false activation issue. Only the proposed method successfully solves the problem and remains a high dice score.
% In the overall evaluation, our AME-CAM achieves the highest dice score and IoU among all the comparisons with other CAM-based approaches. Moreover, the dice score and the IoU are over 0.8 in the T2-FLAIR modality of the BraTS dataset.


% {\small
% \bibliographystyle{splncs04}
% \bibliography{miccai23}
% }

\end{document}
